\section{Eleven positive variables: combining quotient absorption with a smaller norm}
\label{sec:eleven}

Package 38 removes the Pell index quotient $h$ by absorbing its
divisibility condition into an existing carrier. Package 44 removes an
inner radical separator, but its displayed root estimate uses $k>E$,
which was obtained from that same quotient equation. Combining the two
constructions therefore requires a new estimate, not just substituting
one formula into the other. The estimate below compares the coupled-root
gap with the actual relative error. It works before $h$ has been
reconstructed and preserves the smaller weight $K=IT+1$.

\subsection{The two explicit equations}

The supplied variables are precisely
\[
 J,F,\alpha,z,i,T,s,w,\eta,\zeta,t\in\Pos.
\]
Use the complete fixed tuple $(\beta,K_0,M_C,M_F)$ specified in the
coefficient appendix. The following names are polynomial abbreviations,
not additional unknowns:
\begin{align*}
 q&=\beta J+1,&s_0&=s+3q,&C_w&=q-F-\alpha,\\
 X&=wq,&Y&=s_0q^3,&E&=XY,\\
 k&=\eta+\zeta,&c&=kY+\eta,&a&=E+Y,\\
 H&=4a+3,&\Delta&=a^2+H=(a+2)^2-1,&\ell&=EkY,\\
 A&=1+\ell(\ell+k),&G&=1+\Delta c^2,&b_0&=X+ac,\\
 I&=ic^2,&D&=1+\Delta I^2,&Q&=\Delta(D-1)=(\Delta I)^2,\\
 d&=Q-1,&v&=cT.
\end{align*}
Define also
\begin{align*}
 R&=((q-1)(q-F)+\alpha)(q^2-1)+(M_C+qM_F)J,\\
 N_t&=(K_0+w)C_w+(q-F)-z(q-1),\\
 b&=c+RD,\qquad e=Q(Dv^2+b^2)-1,\qquad B_a=2Qvb.
\end{align*}
Choose one of two carrier corrections,
\[
 \delta_{\mathrm f}=(a+2)c-\ell
 \quad\hbox{or}\quad
 \delta_{\mathrm s}=ac-\ell.
\]
The subscripts denote the full and short carriers. For either choice put
\begin{equation}\label{eq:eleven-carrier}
 \begin{split}
 K&=IT+1,\qquad S=EH,\qquad L=E(1+HK),\\
 W&=Eb_0+H\bigl(E(t+b+\delta)+k-R-1\bigr).
 \end{split}
\end{equation}
The computed carrier weight $K$ is distinct from the fixed compiler
coefficient $K_0$.

Define the integer polynomial
\begin{equation}\label{eq:eleven-norm}
 \mathcal N=
 \prod_{\substack{g=\pm\sqrt G\\ \tau=\pm\sqrt A\\ f=\pm\sqrt D}}
 \left[d(W-Lg-S\tau)^2-S^2e+S^2B_af\right].
\end{equation}
The signs in this display specify a formal eight-factor product. Each
root appears through an even polynomial after the product is expanded,
so $\mathcal N$ has integer coefficients. The notation does not add
quantified variables or assert that independent Galois automorphisms of
the three radicals exist. The requested scalar equation is
\begin{equation}\label{eq:eleven-polynomial}
 \boxed{\quad P_{11}=N_t(1+\mathcal N)-1=0.\quad}
\end{equation}

For an entirely division-free evaluation of \eqref{eq:eleven-norm}, set
\begin{align*}
 n_d&=dS^2A,&H_P&=S^2QDv^2,&H_U&=S^2Qb^2,\\
 H_e&=H_P+H_U-S^2,&j_0&=4H_PH_U,\\
 \mathcal T&=2d(W^2+L^2G),&\mathcal K&=[d(W^2-L^2G)]^2,\\
 A_1&=2(n_d+H_e),&c_1&=(n_d-H_e)^2-j_0,\\
 a_1&=\mathcal T-A_1,&b_1&=c_1-\mathcal K,\\
 v_1&=16n_dj_0-a_1(a_1\mathcal T+2b_1),&
 u_1&=b_1^2-a_1^2\mathcal K.
\end{align*}
Then
\begin{equation}\label{eq:eleven-scalar-norm}
 \mathcal N=\mathcal K v_1^2-u_1(\mathcal T v_1-u_1).
\end{equation}
One verifies this identity by eliminating $\tau,f$ first and reducing
the resulting quartic in $d(W-Lg)^2$ modulo its quadratic relation.
It is a polynomial identity, including at square or coincident
radicands, not an identity restricted to successful witnesses.

\begin{theorem}\label{thm:eleven}
For each choice of $\delta$, the polynomial
\eqref{eq:eleven-polynomial} has a zero in $\Pos^{11}$ if and only if
the normalized blank-K215 compiler has a solution. Thus its absence of
positive zeros is provably equivalent in ZFC to $\Con(\ZFC)$, through
the inherited compiler and machine/refutation calibration.
The full carrier has exact degree $990$ and a $117$-operation evaluator;
the short carrier has exact degree $994$ and a $115$-operation evaluator.
\end{theorem}

\subsection{Extracting the unit and obtaining the outer bounds}

Put $x_0=d(W^2+L^2G)$. The scalar evaluator satisfies
\[
 \mathcal T=2x_0,\qquad
 \mathcal K=x_0^2-4d^2W^2L^2G,
\]
and consequently
\begin{equation}\label{eq:eleven-modfour}
 \mathcal N\equiv(u_1-x_0v_1)^2\pmod4.
\end{equation}
At a zero of $P_{11}$ both $N_t$ and $1+\mathcal N$ are integer
units. The possibility $1+\mathcal N=-1$ would give
$\mathcal N=-2$, contradicting \eqref{eq:eleven-modfour}. Therefore
\[
 \mathcal N=0,\qquad N_t=1.
\]
The authentic coefficient interface now gives
\begin{equation}\label{eq:eleven-bounds}
 \begin{gathered}
 q\ge16,\quad C_w>0,\quad 3q<R<q^4-2,\\
 X\ge q,\quad Y=(s+3q)q^3>3q^4>3R,\quad E>R+1,\\
 E\ge65536,\quad k\ge2,\quad kY<c<k(Y+1)<2kY,\\
 I\ge c^2,\quad c\ge2,\quad a+2<2E.
 \end{gathered}
\end{equation}
For example, $C_w\le0$ would give
$N_t\le1-F-(z-1)(q-1)\le0$, so the boundary $C_w=0$ is excluded
here by the already recovered transport sign. No assertion $k>E$,
$c>EY$, or $R<k$ is used before the omitted quotient is recovered.

\subsection{A relative-error estimate replaces the lost quotient bound}

For estimates only, let $g=\sqrt G$, $f=\sqrt D$ and $\tau=\sqrt A$
denote the positive real roots. From \eqref{eq:eleven-bounds},
\begin{equation}\label{eq:eleven-roots}
 \begin{gathered}
 g<(a+2)c<2Ec,\qquad f<(a+2)I<2EI,\\
 1<A<\Delta I=\sqrt Q,\qquad g>\tau.
 \end{gathered}
\end{equation}
For the radicand comparison, $\Delta c^2-A$ exceeds
$k^2[2EY^3+Y^4-EY]-1>0$; also $I\ge c^2$. These estimates do
not require the index equation.

\begin{lemma}[The sufficient coupled-root gap]\label{lem:eleven-gap}
Under \eqref{eq:eleven-bounds},
\[
 \frac{b+cf}{d}<Ig-cf.
\]
Every coherent auxiliary root, including a complex root, satisfying
\[
 d y^2=Q(vf'-b)^2-1=e-B_af',\qquad f'=\pm f,
\]
therefore obeys
\[
 |y|<b+ITg,\qquad |\tau'+y|<b+(IT+1)g
 \quad(\tau'=\pm\tau).
\]
\end{lemma}

\begin{proof}
The exact identity $(Ig)^2-(cf)^2=I^2-c^2$ first gives $Ig>cf$.
Since $I\ge c^2$ and $c\ge2$, its right side is greater than
$I^2/2$. Thus
\[
 Ig-cf>\frac{I^2}{4Ig}>\frac{I}{8Ec}
       \ge\frac{c}{8E}>\frac{Y}{4E}.
\]
On the other hand $c<I^2$, $D<4E^2I^2$, $R<E$ and $cf<2EI^2$
give
\[
 b<5RE^2I^2<5E^3I^2,\qquad
 b+cf<6E^3I^2,\qquad d>\tfrac12E^4I^2.
\]
Consequently
\[
 \frac{b+cf}{d}<\frac{12}{E}<\frac{Y}{4E}<Ig-cf,
\]
because $Y>48$. The difference $Ig-cf$ need not exceed one: comparison
with this smaller error is the new point.

Set $Z=b+Tcf>1$. Taking complex absolute values and using $Q=d+1$,
\[
 |y|^2\le\frac{QZ^2+1}{d}
 <\bigl[Z(1+1/d)\bigr]^2.
\]
The last strict difference is $(QZ^2-d)/d^2>0$. Since $T\ge1$,
\[
 Z/d\le T(b+cf)/d<T(Ig-cf).
\]
Hence $|y|<Z+Z/d<b+ITg$. Finally $\tau<g$ proves the stated
bound for $|\tau'+y|$.
\end{proof}

\subsection{Recovering the radicals and the absorbed divisibility}

Eliminating the two inner radicals gives the rational-coefficient
polynomial
\begin{align*}
 C(x)&=d(x^2+A)-e,\\
 \mathcal F(x)&=[C(x)^2+DB_a^2-4d^2x^2A]^2-4C(x)^2DB_a^2.
\end{align*}
At $\mathcal N=0$ some branch gives
$x=(W-Lg')/S=\tau'+y$ and $\mathcal F(x)=0$. If $g'=-g$, then
\[
 x=b+\left[t+\frac{k-R-1}{E}\right]
      +\delta+\frac{b_0}{H}+(K+1/H)g>b+Kg.
\]
The bracket is positive because $t\ge1$, $k\ge2$ and $E>R+1$;
both carrier corrections are positive. This contradicts
Lemma~\ref{lem:eleven-gap}. If $G$ were nonsquare, the actual quadratic
conjugation in $\Q(\sqrt G)$ would carry a root at one main sign to
the excluded sign. Therefore $g$ is a positive integer and $x$ is
rational.

The exceptional inner branches can be excluded without a separator.
If $C(x)=0$ at a root, substitution gives the integer identity
\[
 d^2A^2-deA+Q^2v^2b^2D=0.
\]
Modulo $Q$, both $d$ and $e$ are $-1$, so $Q\mid A(A-1)$.
This is impossible because $1<A<\sqrt Q$. Likewise
$\mathcal F(0)=0$ would force $Q\mid(A-1)^2$, also impossible.
Thus both denominators in the following reconstruction are nonzero:
\begin{align*}
 U_x&=C(x)^2+DB_a^2-4d^2x^2A,\\
 f'&=-\frac{U_x}{2C(x)B_a},\qquad
 \tau'=\frac{C(x)+B_af'}{2dx},\qquad y=x-\tau'.
\end{align*}
The equation $\mathcal F(x)=0$ gives $(f')^2=D$, $(\tau')^2=A$
and the auxiliary equation. The rational roots $f',\tau'$ are
integers.

We use the inherited rational Pell saturation lemma: for integers
$m\ge2,V$, a rational solution of
$m^2V^2-(m^2-1)y^2=1$ is integral. Briefly, writing $y=u/v$ in
lowest terms gives $v^2\mid m^2-1$. In the integer-coordinate Pell
group of discriminant $(m^2-1)/v^2$, the solutions $(m,v)$ and
$(m|V|,|u|)$ have indices $j,l$. Divisibility of their abscissas
forces $l$ to be an odd multiple of $j$: reduce the index modulo $2j$,
reflect into $[0,j]$, and use strict growth of the abscissas. The
ordinate divisibility then gives $v\mid u$, hence $v=1$.
Apply this lemma with $m=\Delta I$ and $V=vf'-b$. It makes $y$
and consequently $x$ integral. Reducing $W=Lg+Sx$ modulo $E$ yields
\begin{equation}\label{eq:eleven-weak-index}
 E\mid H(k-R-1).
\end{equation}
The factor $H$ has not yet been cancelled.

\subsection{The short-range argument restores the full index}

Let $\chi_{a+2},\psi_{a+2}$ be the Pell sequences. The recovered
main and normalized strong equations give indices $p,m\ge1$ with
\[
 g=\chi_{a+2}(p),\quad c=\psi_{a+2}(p),\quad
 |f'|=\chi_{a+2}(m),\quad I=\psi_{a+2}(m).
\]
Since $c^2\mid I$, strong divisibility and composition first give
$p\mid m$ and then $c\mid m/p$. In particular $c\mid m$ and
$m\ge pc>2p$. The auxiliary odd-index identities give a further
index $j_0$ such that
\[
 j_0\equiv\pm p\pmod m,\qquad j_0\equiv\pm R\pmod c.
\]
These are the normalized signed-quotient identities used in
package 38: they follow from $V\equiv-c\pmod{|f'|}$,
$V\equiv-R\pmod c$, the odd Pell quotient polynomial, and signed
index reduction. They allow either sign of $f'$ or $V$. As
$c>2Y>6R$ and $c>2p$, they imply $p=R$ before any use of $h$.
For the latter bound, $c>1$ gives $p\ge2$, and the Pell recurrence
gives $c\ge5^{p-1}>2p$.

The first equation classifies
\[
 |\tau'|=\chi_{2EY+1}(n),\qquad
 k=2\psi_{2EY+1}(n),\qquad 1\le n<p=R.
\]
The last inequality follows by parameter monotonicity from
$2EY+1>a+2$ and $c>k$. Put $\mu=E/\gcd(E,H)$. Because
$\gcd(Y,H)=\gcd(Y,3)\le3$,
\[
 \mu\ge Y/3>R.
\]
Equation \eqref{eq:eleven-weak-index} gives $k\equiv R+1\pmod\mu$.
The first Pell recurrence gives $k\equiv2n\pmod E$. Hence
$2n\equiv R+1\pmod\mu$; since $|2n-R-1|<R<\mu$, this means
$2n=R+1$ exactly. It follows that
\[
 h=(k-R-1)/E\in\Pos.
\]
Indeed $R>48$ gives $n\ge25$, so strict Pell growth gives $k>2n$.

Now divide the carrier identity by $E$ and reduce modulo $H$. This
gives $g\equiv b_0\pmod H$, so $\gamma=(g-b_0)/H$ is an integer.
It is positive: $n\ge25$ implies $k\ge4(2EY+1)>X$, hence $c>X$,
whereas $g>(a+1)c$. We have recovered every equation of the normalized
compiler. Taking absolute values of $\tau',y$ is harmless; the
inherited signed-quotient interface permits the recovered auxiliary
signs. This proves soundness using the identified compiler theorem,
not by assuming its conclusion during radical recovery.

\subsection{Completeness and exact complexity}

For a blank halting run, choose canonical positive normalized compiler
witnesses. Their estimates give $Y>2^R$, $R>3q$ and $q\ge16$, so
$Y/q^3>3q$ and $s=Y/q^3-3q$ is positive. For both corrections,
$0<\delta<g$: since $c<2kY$, one has $\ell/c>E/2>2$, and
$(a+2)c-\ell<ac<g$. Define
\[
 t=\gamma+Kg+\tau+y-b-\delta-h.
\]
The auxiliary equation gives $y\ge|Tcf-b|\ge b-Tcf$.
Consequently
\[
 t>\gamma+\tau-h+T(Ig-cf)>0,
\]
because $\tau>\ell>k/E>h$ and $Ig>cf$. Its defining identity gives
$W=Lg+S(\tau+y)$, so a coherent factor in
\eqref{eq:eleven-norm} vanishes and $N_t=1$. This proves completeness.

The exact degrees can be transferred without assuming generic
coefficients. Let $\mathcal N_{44}$ be package 44's corresponding
full or short norm, before quotient absorption. In the localization at
$E$, there is the exact identity
\begin{equation}\label{eq:eleven-degree-transfer}
 \mathcal N=
 E^{16}\mathcal N_{44}\left(s+3q,\,
                      t+\frac{k-R-1}{E}\right).
\end{equation}
Each of its eight factors is multiplied by $E^2$. The old norm has
$t$-degree at most $16$, so denominators clear and this is an ordinary
polynomial identity. Its source exact degrees are $892$ and $896$
for the full and short carriers respectively \cite{package44}.
The rational shift has degree $4-6=-2$, strictly below $\deg t=1$;
any term using it loses at least three degrees. The highest part of
$s\mapsto s+3q$ is the invertible linear substitution
$s\mapsto s+3\beta J$, so it preserves the nonzero leading form.
Since $\deg E=6$, \eqref{eq:eleven-degree-transfer} gives exact
norm degrees $988$ and $992$. Finally $N_t$ has nonzero leading form
$w(\beta J-F-\alpha)-z\beta J$, of degree two, giving exact scalar
degrees $990$ and $994$.

The literal division-free DAGs charge all operations in these definitions
and contain no dead gates. They give $117$ operations for the full
carrier and $115$ for the short carrier. The degree proof, the charged
gate counts, and the semantic equivalence are separate claims: the
finite circuit audits support the first two, while the argument above
and its explicit inherited Pell/compiler dependencies establish the
third. No new end-to-end formal certificate or absolute minimality
claim follows from this construction.
